% Extracción quirúrgica del propietario V2 05d: líneas 286--512.
% El resto permanece en este archivo como procedencia, pero no se publica.
\iffalse
% Macroestructura de coherencia estrictamente interna del continuo discreto.
% No usa geometrías, clases ni conclusiones clásicas.

\chapter{Macrostructure interne de cellules, cohérence et complétude}
\label{chap:pdfv2-coh-interna}

\begin{statusbox}
L'unique entrée de ce chapitre est la structure discrète du continu. La
branche \(\mathrm{coh}\) se construit sur les anneaux
\(A_N=\mathbb Z/3^N\mathbb Z\), les cellules APP, leurs histoires et leurs
terminaux. Sa chaîne non ablative est
\[
 \mathcal C_{\rm cont}^{\rm disc}
 \supset\mathcal C_{\rm cont}^{\rm coh}
 \xrightarrow{\operatorname{Res}_{\rm coh}}
 \mathcal G_{\rm coh}
 \xrightarrow{\operatorname{pr}_{X,{\rm coh}}}
 X_{\chi,{\rm coh}}
 \xrightarrow{\mathcal A_{\rm coh}}
 \mathfrak M_{\rm coh}
 \xrightarrow{\mathcal P_{\rm coh}}
 \mathcal C_{\rm coh}^{\rm HMT}.
\]
Aucun objet classique n'apparaît en entrée ni en sortie. Chaque image
conserve littéralement son histoire génératrice.
\end{statusbox}

\section{Domaine typé et séparation des indices}
\label{sec:pdfv2-coh-domain}

Pour ne pas identifier profondeur et précision, on utilise l'ensemble dirigé
\begin{equation}\label{eq:pdfv2-coh-lambda}
 \Lambda:=\mathbb N_{\rm prof}\times\mathbb N_{\rm prec},
 \qquad
 \lambda=(d,N),
 \qquad
 A_N:=\mathbb Z/3^N\mathbb Z.
\end{equation}
Si \(\lambda\leq\lambda'=(d',N')\), la transition
\begin{equation}\label{eq:pdfv2-coh-rho-lambda}
 \rho_{\lambda',\lambda}:
 \mathcal C_{{\rm cont},\lambda'}^{\rm disc}
 \longrightarrow
 \mathcal C_{{\rm cont},\lambda}^{\rm disc}
\end{equation}
tronque l'histoire après la profondeur \(d\) et réduit les coefficients par
\(A_{N'}\twoheadrightarrow A_N\), en conservant feuillet, orientation, résidu,
quotient, retenue, mémoire, frontière, chemin, survie et terminal.

La sous-structure \(\mathcal C_{\rm cont}^{\rm coh}\) se compose des histoires
\(x\in\mathcal C_{\rm cont}^{\rm disc}\) qui satisfont, à tout niveau :
\begin{enumerate}
 \item le préfixe de \(x\) porte les cellules orientées et basées définies dans
       la section suivante ;
 \item chaque raffinement induit une application du complexe de la cellule ;
 \item troncature, réduction des coefficients, orientation et extension
       nonadique commutent ;
 \item les multisections d'extension survivante sont non vides à toute
       profondeur ;
 \item le retour nonadique conserve la différence terminale entre l'histoire
       et sa continuation.
\end{enumerate}
Ces conditions définissent le domaine de la branche. Elles ne sont pas attribuées à une
histoire arbitraire sans être vérifiées.

\section{Cellule APP sur \texorpdfstring{\(A_N\)}{AN}}
\label{sec:pdfv2-coh-cell}

Une cellule orientée est
\[
 \nu=(h_\nu,I_\nu,J_\nu,\epsilon_\nu),
\]
où \(h_\nu\) est une histoire finie, \(I_\nu,J_\nu\) sont des ensembles
ordonnés non vides de ports,
\[
 a_\nu:=|I_\nu|\geq1,
 \qquad
 b_\nu:=|J_\nu|\geq1,
\]
et \(\epsilon_\nu\in\{+1,-1\}\) est son orientation.

On définit
\begin{equation}\label{eq:pdfv2-coh-kappa}
 \begin{aligned}
 \kappa_{\nu,N}:
 A_N^{I_\nu}\oplus A_N^{J_\nu}
 &\longrightarrow A_N^{I_\nu\times J_\nu},\\
 \kappa_{\nu,N}(u,v)_{ij}&:=u_i-v_j.
 \end{aligned}
\end{equation}
Les ports de base \(i_0\in I_\nu\), \(j_0\in J_\nu\) étant choisis, soit
\begin{equation}\label{eq:pdfv2-coh-delta}
 \begin{aligned}
 \delta_{\nu,N}:A_N^{I_\nu\times J_\nu}
 &\longrightarrow
 A_N^{(I_\nu\setminus i_0)\times(J_\nu\setminus j_0)},\\
 \delta_{\nu,N}(w)_{ij}
 &:={}
 w_{ij}-w_{i j_0}-w_{i_0j}+w_{i_0j_0}.
 \end{aligned}
\end{equation}

\begin{proposition}[Suite exacte de la cellule]
\label{prop:pdfv2-coh-cell-exact}
Pour toute cellule \(\nu\) et toute précision \(N\), la suite
\begin{equation}\label{eq:pdfv2-coh-exact-sequence}
 0\longrightarrow A_N
 \xrightarrow{\ \Delta\ }
 A_N^{I_\nu}\oplus A_N^{J_\nu}
 \xrightarrow{\ \kappa_{\nu,N}\ }
 A_N^{I_\nu\times J_\nu}
 \xrightarrow{\ \delta_{\nu,N}\ }
 A_N^{(a_\nu-1)(b_\nu-1)}
 \longrightarrow0,
\end{equation}
où
\[
 \Delta(t)=\bigl((t)_{i\in I_\nu},(t)_{j\in J_\nu}\bigr),
\]
est exacte et scindée sur \(A_N\).
\end{proposition}

\begin{proof}
Si \(\kappa_{\nu,N}(u,v)=0\), alors \(u_i=v_j\) pour tout \(i,j\),
de sorte que \((u,v)\) est diagonale. Ainsi,
\(\ker\kappa_{\nu,N}=\operatorname{im}\Delta\). La substitution directe donne
\(\delta_{\nu,N}\kappa_{\nu,N}=0\).

Si \(\delta_{\nu,N}(w)=0\), posons
\[
 u_i=w_{i j_0},
 \qquad
 v_j=w_{i_0j_0}-w_{i_0j}.
\]
Alors
\[
 u_i-v_j=w_{i j_0}+w_{i_0j}-w_{i_0j_0}=w_{ij},
\]
de sorte que \(\ker\delta_{\nu,N}=\operatorname{im}\kappa_{\nu,N}\). Pour
prouver la surjectivité de \(\delta_{\nu,N}\), on annule la ligne et la
colonne de base et l'on place arbitrairement les coordonnées souhaitées aux
positions restantes. Ces mêmes formules fournissent des scindements.
\end{proof}

La cellule complexe, aux degrés \(1\to0\), est le cône homologique de
\(\kappa_{\nu,N}\) avec la convention qui place sa source au degré \(1\) :
\begin{equation}\label{eq:pdfv2-coh-cell-complex}
 C_{\nu,N}:=\operatorname{Cone}(\kappa_{\nu,N})
 :=\left[
 A_N^{I_\nu}\oplus A_N^{J_\nu}
 \xrightarrow{\ \kappa_{\nu,N}\ }
 A_N^{I_\nu\times J_\nu}
 \right].
\end{equation}
La proposition précédente donne
\begin{equation}\label{eq:pdfv2-coh-cell-homology}
 H_1(C_{\nu,N})\simeq A_N,
 \qquad
 H_0(C_{\nu,N})\simeq A_N^{(a_\nu-1)(b_\nu-1)}.
\end{equation}
Le défaut de cellule est ainsi construit, non ajouté comme un cardinal sans
application.

\section{Orientation et retenue de la cellule}
\label{sec:pdfv2-coh-orientation-carry}

Soient
\[
 \tau_1(u,v)=(v,u),
 \qquad
 P(w)_{ji}=w_{ij}.
\]
L'échange orienté est l'application de complexes
\begin{equation}\label{eq:pdfv2-coh-theta}
 \Theta_{\nu,N}:=(\tau_1,-P):
 C_{a_\nu,b_\nu,N}\longrightarrow C_{b_\nu,a_\nu,N},
\end{equation}
parce que
\begin{equation}\label{eq:pdfv2-coh-orientation}
 \kappa_{b_\nu,a_\nu,N}\tau_1
 =-P\kappa_{a_\nu,b_\nu,N}.
\end{equation}
En outre,
\(\Theta_{\nu^{\rm op},N}\Theta_{\nu,N}=I\). Le chemin opposé ne
s'identifie pas à l'original : il change le signe de degré zéro et conserve l'acte
d'inversion.

Pour \(r\in A_N\), soit
\(\langle r\rangle_N\in\{0,\ldots,3^N-1\}\) son représentant canonique.
L'emprunt/retenue local est
\begin{equation}\label{eq:pdfv2-coh-beta}
 \beta_N(r,s):=
 \frac{\langle r-s\rangle_N-\langle r\rangle_N+\langle s\rangle_N}
 {3^N}\in\{0,1\}.
\end{equation}
Ainsi,
\begin{equation}\label{eq:pdfv2-coh-carry-cell}
 \langle u_i-v_j\rangle_N
 =\langle u_i\rangle_N-\langle v_j\rangle_N
  +3^N\beta_N(u_i,v_j).
\end{equation}
La matrice
\[
 \operatorname{Carr}_{\nu,N}(u,v)
 :=\bigl(\beta_N(u_i,v_j)\bigr)_{ij}
\]
est conservé avec \(\kappa_{\nu,N}(u,v)\) : résidu et retenue sont deux
coordonnées distinctes.

Lorsque les tailles des ports proviennent des deux feuillets APP, l'admissibilité
de la cellule inclut l'identité entière
\begin{equation}\label{eq:pdfv2-coh-app-defect}
 (a_\nu-1)(b_\nu-1)
 =\rho_{\Pi,\nu}-\rho_{\Sigma,\nu}+1
  +9(c_{\Pi,\nu}-c_{\Sigma,\nu}).
\end{equation}
L'égalité est exigée uniquement pour une cellule ayant cette généalogie ; elle n'est pas affirmée
pour des entiers dépourvus de feuillets APP.

Si \(L_N(u)\) est le relèvement entier canonique de la matrice d'une flèche
\(u\), on définit la retenue de composition par
\begin{equation}\label{eq:pdfv2-coh-carry-composition}
 L_N(v)L_N(u)=L_N(vu)+3^N c_N(v,u).
\end{equation}
L'associativité de trois flèches donne le cocycle exact
\begin{equation}\label{eq:pdfv2-coh-carry-cocycle}
 c_N(w,vu)+L_N(w)c_N(v,u)
 =c_N(wv,u)+c_N(w,v)L_N(u).
\end{equation}

\section{Catégorie d'histoires et foncteur cellulaire}
\label{sec:pdfv2-coh-category}

Pour un préfixe \(x_{\leq d}\), soit \(J_d(x)\) la catégorie finie dont les
objets sont toutes ses cellules \(\nu\). Ses flèches sont engendrées par :
\begin{enumerate}
 \item des raffinements \(u:\nu\to\nu'\) avec des applications surjectives de
       ports
       \[
        \alpha_u:I_{\nu'}\twoheadrightarrow I_\nu,
        \qquad
        \beta_u:J_{\nu'}\twoheadrightarrow J_\nu;
       \]
 \item les échanges d'orientation \(\Theta_\nu\) ;
 \item les extensions nonadiques \(\Gamma_9\nu\) ;
 \item l'identité et la composition, soumises aux relations de frontière,
       de chemin, d'orientation et de retenue de l'histoire.
\end{enumerate}
La finitude se rapporte au préfixe ; le système de tous les préfixes n'est pas
déclaré fini.

Le foncteur cellulaire est
\begin{equation}\label{eq:pdfv2-coh-gamma-functor}
 \Gamma_{d,N}:J_d(x)\longrightarrow
 \operatorname{Ch}^{[0,1]}(A_N),
 \qquad
 \Gamma_{d,N}(\nu)=C_{\nu,N}.
\end{equation}
Pour une flèche de raffinement \(u:\nu\to\nu'\), il agit par
\begin{align}
 \Gamma_{d,N}(u)_1(p,q)
 &:=(p\circ\alpha_u,q\circ\beta_u),
 \label{eq:pdfv2-coh-gamma-one}\\
 \Gamma_{d,N}(u)_0(w)
 &:=w\circ(\alpha_u\times\beta_u).
 \label{eq:pdfv2-coh-gamma-zero}
\end{align}
On vérifie coordonnée par coordonnée
\begin{equation}\label{eq:pdfv2-coh-gamma-chain-map}
 \kappa_{\nu',N}\Gamma_{d,N}(u)_1
 =\Gamma_{d,N}(u)_0\kappa_{\nu,N}.
\end{equation}
Sur une inversion d'orientation, il agit par
\(\Theta_{\nu,N}\), et sur les compositions, il agit par composition, en gardant
séparément \(c_N(v,u)\).

La réduction \(r_{N+1,N}:A_{N+1}\twoheadrightarrow A_N\) satisfait
\begin{equation}\label{eq:pdfv2-coh-kappa-reduction}
 r_{N+1,N}^{I\times J}\kappa_{\nu,N+1}
 =\kappa_{\nu,N}
  \bigl(r_{N+1,N}^{I}\oplus r_{N+1,N}^{J}\bigr).
\end{equation}
Par conséquent, le diagramme des cellules est naturel simultanément en histoire et en
précision.

\fi
\section{Module des chemins, totalisation cellulaire et terminal nonadique}
\label{sec:pdfv2-coh-route-module}

Soit \(\operatorname{Term}_d(x)\) l'ensemble des terminaux de mémoire du
préfixe. Pour \(\nu\in J_d(x)\), définissons
\begin{equation}\label{eq:pdfv2-coh-w-module}
 W_{d,N}(\nu):=
 A_N\bigl[
 \operatorname{Path}_{J_d(x)}(\nu,\operatorname{Term}_d(x))
 \bigr].
\end{equation}
C'est le module libre sur tous les chemins terminaux ; il n'en sélectionne aucun. Si
\(u:\nu\to\nu'\), l'action contravariante est
\begin{equation}\label{eq:pdfv2-coh-w-action}
 W_{d,N}(u):W_{d,N}(\nu')\longrightarrow W_{d,N}(\nu),
 \qquad
 [\lambda]\longmapsto[\lambda\circ u].
\end{equation}
Chaque générateur conserve orientation, retenue, frontière, chemin et mémoire.

\section{Complexe bar bilatéral et totalisation de la bichaîne}
\label{sec:pdfv2-coh-bar}

Pour \(q\geq0\) et \(r\in\{0,1\}\), soit
\begin{equation}\label{eq:pdfv2-coh-bar-bigraded}
 B_{q,r}^{d,N}(x):=
 \bigoplus_{\nu_0\xrightarrow{u_1}\cdots\xrightarrow{u_q}\nu_q}
 W_{d,N}(\nu_q)\otimes_{A_N}\Gamma_{d,N}(\nu_0)_r.
\end{equation}
La différentielle bar est définie sans abréviation par
\begin{align}
 d_{\rm bar}(w;u_1,\ldots,u_q;c)
 :={}&(w;u_2,\ldots,u_q;\Gamma(u_1)c)
 \nonumber\\
 &+\sum_{i=1}^{q-1}(-1)^i
 (w;u_1,\ldots,u_{i+1}u_i,\ldots,u_q;c)
 \nonumber\\
 &+(-1)^q
 (W(u_q)w;u_1,\ldots,u_{q-1};c).
 \label{eq:pdfv2-coh-bar-differential}
\end{align}
Pour \(q=1\), c'est exactement la relation
\begin{equation}\label{eq:pdfv2-coh-coend-relation}
 d_{\rm bar}(w;u;c)
 =w\otimes\Gamma(u)c-W(u)w\otimes c.
\end{equation}
On identifie ainsi une incidence transportée à la même incidence écrite
dans la carte raffinée ; la relation ne reste pas cachée sous le nom du
totalisateur.

Le complexe total est
\begin{equation}\label{eq:pdfv2-coh-total-complex}
 K_{d,N}(x):=\operatorname{Tot}B_{\bullet,\bullet}^{d,N}(x),
 \qquad
 D_{\rm tot}:=d_{\rm bar}+(-1)^q\partial_\Gamma
 \quad\text{sur }B_{q,r}^{d,N}.
\end{equation}
La fonctorialité de \(W\) et \(\Gamma\), avec
\(\partial_\Gamma^2=0\), donne
\begin{equation}\label{eq:pdfv2-coh-total-square-zero}
 D_{\rm tot}^2=0.
\end{equation}
En notation compacte, mais désormais établie par
\eqref{eq:pdfv2-coh-bar-bigraded}--
\eqref{eq:pdfv2-coh-bar-differential},
\begin{equation}\label{eq:pdfv2-coh-derived-coend}
 K_{d,N}(x)
 \simeq
 W_{d,N}\otimes^{\mathbf L}_{A_N[J_d(x)]}\Gamma_{d,N}.
\end{equation}

\section{Colimite de profondeur et terminal nonadique}
\label{sec:pdfv2-coh-terminal}

Les préfixes forment la colimite
\begin{equation}\label{eq:pdfv2-coh-k-infinity}
 K_{\infty,N}(x):=\varinjlim_d K_{d,N}(x).
\end{equation}
L'extension de neuf phases fait passer de la profondeur \(d\) à la profondeur \(d+9\).
Ce n'est qu'après la formation de \eqref{eq:pdfv2-coh-k-infinity} qu'elle induit un
endomorphisme bien typé
\begin{equation}\label{eq:pdfv2-coh-u-gamma9}
 U_{\Gamma_9,N}:K_{\infty,N}(x)
 \longrightarrow K_{\infty,N}(x).
\end{equation}
Le terminal est
\begin{equation}\label{eq:pdfv2-coh-terminal-cone}
 \mathcal T_{{\rm coh},N}(x)
 :=\operatorname{Cone}\bigl(I-U_{\Gamma_9,N}\bigr).
\end{equation}
Avec la convention homologique
\[
 \operatorname{Cone}(f)_q
 =K_{\infty,N,q-1}\oplus K_{\infty,N,q},
\]
sa différentielle est
\begin{equation}\label{eq:pdfv2-coh-cone-differential}
 \partial_{\rm Cone}(a,b)
 =\bigl(-\partial_Ka,\partial_Kb+f(a)\bigr).
\end{equation}

Bien que la phase visible de \(h\) et \(\Gamma_9h\) coïncide, les générateurs
\([h]\) et \([\Gamma_9h]\) sont différents parce que leurs ledgers et leurs profondeurs
sont différents. Le terminal conserve
\begin{equation}\label{eq:pdfv2-coh-terminal-difference}
 (I-U_{\Gamma_9,N})[h]=[h]-[\Gamma_9h].
\end{equation}
Les réductions de précision commutent avec le terminal :
\begin{equation}\label{eq:pdfv2-coh-u-reduction}
 r_{N+1,N}U_{\Gamma_9,N+1}
 =U_{\Gamma_9,N}r_{N+1,N}.
\end{equation}

\section{Lecteur cohomologique terminal : cycle d'incidence et classe dans \(H_0(\operatorname{Cone}(I-U_{\Gamma_9,N}))\)}
\label{sec:pdfv2-coh-reader}

Chaque cellule de \(x_{\leq d}\) porte un vecteur d'incidence
\[
 \omega_{\nu,N}(x)
 \in A_N^{I_\nu\times J_\nu}=C_{\nu,N,0}
\]
et le chemin terminal complet \(\lambda_{\nu,x}\). Le cycle brut est
\begin{equation}\label{eq:pdfv2-coh-z-cycle}
 z_{d,N}(x):=
 \sum_{\nu\in\operatorname{Cell}_d(x)}
 \epsilon_\nu[\lambda_{\nu,x}]
 \otimes\omega_{\nu,N}(x)
 \in B_{0,0}^{d,N}(x).
\end{equation}
Comme le complexe total est homologique non négatif,
\begin{equation}\label{eq:pdfv2-coh-z-closed}
 D_{\rm tot}z_{d,N}(x)=0.
\end{equation}
On conserve aussi bien le cycle que sa classe
\[
 [z_{d,N}(x)]\in H_0(K_{d,N}(x)).
\]
Après son inclusion dans \(K_{\infty,N}\), sa classe terminale est
\[
 [(0,z_{d,N}(x))]
 \in H_0(\mathcal T_{{\rm coh},N}(x)).
\]
Le lecteur complet est
\begin{equation}\label{eq:pdfv2-coh-reader-full}
 \mathcal L_{{\rm coh},d,N}(x)
 :=\bigl(
 z_{d,N}(x),[z_{d,N}(x)],
 U_{\Gamma_9,N}z_{d,N}(x),
 [(0,z_{d,N}(x))]
 \bigr).
\end{equation}
On ne publie pas uniquement une classe quotient : le représentant, le
chemin qui l'a produit, sa continuation et le terminal subsistent.

\section{Fibres survivantes et corrections sans sélecteur}
\label{sec:pdfv2-coh-surviving-fibers}

Soit \(\mathscr S_\lambda^{\rm coh}\) le \(A_N\)-module libre engendré par les
histoires finies admissibles au niveau \(\lambda=(d,N)\). Le lecteur se
prolonge linéairement à
\begin{equation}\label{eq:pdfv2-coh-e-lambda}
 e_\lambda:\mathscr S_\lambda^{\rm coh}
 \longrightarrow\mathscr O_\lambda^{\rm coh},
 \qquad
 \mathscr O_\lambda^{\rm coh}
 :=H_0(K_{d,N})\oplus H_0(\mathcal T_{{\rm coh},N}).
\end{equation}
Le codomaine survivant est
\begin{equation}\label{eq:pdfv2-coh-surviving-output}
 \mathscr O_\infty^{\rm surv}
 :=\operatorname{Im}\!\left(
 e_\infty:\mathcal C_{\rm cont}^{\rm coh}
 \longrightarrow\varprojlim_\lambda\mathscr O_\lambda^{\rm coh}
 \right).
\end{equation}
Pour \(a=(a_\lambda)_\lambda\in\mathscr O_\infty^{\rm surv}\), on définit
\begin{equation}\label{eq:pdfv2-coh-fiber}
 \mathscr F_\lambda(a):=
 \left\{s_\lambda\in\mathscr S_\lambda^{\rm coh}:
 \begin{array}{l}
 e_\lambda(s_\lambda)=a_\lambda,\\
 \exists y\in\mathcal C_{\rm cont}^{\rm coh}\text{ avec }
 y|_\lambda=s_\lambda\text{ et }e_\infty(y)=a
 \end{array}
 \right\}.
\end{equation}
Ces fibres sont non vides par la définition du codomaine survivant. Si
\(\lambda'\geq\lambda\), la restriction est surjective :
\begin{equation}\label{eq:pdfv2-coh-fiber-surjection}
 \rho_{\lambda',\lambda}:\mathscr F_{\lambda'}(a)
 \twoheadrightarrow\mathscr F_\lambda(a).
\end{equation}
En effet, une histoire complète qui témoigne de
\(s_\lambda\in\mathscr F_\lambda(a)\) fournit par troncature un
antécédent à tout niveau ultérieur. Cet antécédent n'est pas choisi comme
partie de l'objet.

La multisection complète est
\begin{equation}\label{eq:pdfv2-coh-extension-multisection}
 \Sigma_{\lambda',\lambda}^{\rm coh}(s_\lambda;a)
 :=\{s_{\lambda'}\in\mathscr F_{\lambda'}(a):
 \rho_{\lambda',\lambda}s_{\lambda'}=s_\lambda\}.
\end{equation}
Pour un prolongement provisoire \(\widetilde s_{\lambda'}\), soit
\begin{equation}\label{eq:pdfv2-coh-discrepancy}
 \delta_{\lambda'}:=
 a_{\lambda'}-e_{\lambda'}(\widetilde s_{\lambda'}).
\end{equation}
La relation de toutes les corrections est
\begin{equation}\label{eq:pdfv2-coh-correction-relation}
 \operatorname{Corr}_{\lambda',\lambda}
 (\widetilde s_{\lambda'},\delta_{\lambda'})
 :=\left\{\eta\in\mathscr S_{\lambda'}^{\rm coh}:
 \begin{array}{l}
 \rho_{\lambda',\lambda}\eta=0,\\
 e_{\lambda'}(\eta)=\delta_{\lambda'},\\
 \eta\text{ est supportée sur de nouvelles cellules},\\
 \eta\text{ respecte l’orientation, la retenue, la mémoire et le terminal}
 \end{array}
 \right\}.
\end{equation}
Si deux relèvements sont comparables, leur différence
\(\eta=s_{\lambda'}-\widetilde s_{\lambda'}\) appartient à cette relation
exactement lorsqu'elle satisfait les quatre conditions précédentes. On n'affirme pas
que \eqref{eq:pdfv2-coh-correction-relation} soit non vide pour un
écart arbitraire : la non-vacuité déterminante est celle de
\eqref{eq:pdfv2-coh-extension-multisection} sur le domaine survivant.

\iffalse
\section{Les treize coordonnées de la restriction}
\label{sec:pdfv2-coh-thirteen}

Pour \(x_\lambda\in\mathcal C_{{\rm cont},\lambda}^{\rm coh}\), soit
\begin{equation}\label{eq:pdfv2-coh-d-package}
 \begin{aligned}
 D_{{\rm coh},\lambda}(x_\lambda):=\bigl(&
 \mathcal F_{{\rm coh},\lambda},
 \operatorname{Ext}_{\Gamma,{\rm coh},\lambda},
 \operatorname{Rel}_{{\rm coh},\lambda},
 \mathcal N_{{\rm coh},\lambda},
 \operatorname{or}_{{\rm coh},\lambda},
 \operatorname{Carr}_{{\rm coh},\lambda},
 \operatorname{Mem}_{{\rm coh},\lambda},\\
 &\partial_{{\rm coh},\lambda},
 \gamma_{{\rm coh},\lambda},
 \operatorname{Msect}_{{\rm coh},\lambda},
 \operatorname{Surv}_{{\rm coh},\lambda},
 \operatorname{Term}_{{\rm coh},\lambda},
 \mathcal L_{{\rm coh},\lambda}\bigr).
 \end{aligned}
\end{equation}
Ses coordonnées sont :
\begin{enumerate}
 \item \(\mathcal F_{{\rm coh},\lambda}\) : tous les modules
       \(A_N^{I_\nu}\), \(A_N^{J_\nu}\),
       \(A_N^{I_\nu\times J_\nu}\), les complexes \(C_{\nu,N}\), les
       modules \(W_{d,N}(\nu)\) et \(\mathscr S_\lambda^{\rm coh}\).
 \item \(\operatorname{Ext}_{\Gamma,{\rm coh},\lambda}\) : toutes les applications
       \(\Gamma(u)\), les réductions de coefficients, les troncatures, les extensions
       \(\Gamma_9\) et les relations \(\Sigma_{\lambda',\lambda}\).
 \item \(\operatorname{Rel}_{{\rm coh},\lambda}\) :
       \(\kappa\), \(\delta\), la suite exacte, les faces bar, la
       relation de cofin, l'orientation et les lois de composition.
 \item \(\mathcal N_{{\rm coh},\lambda}\):
       \[
       (d,N,3^N,a_\nu,b_\nu,(a_\nu-1)(b_\nu-1),
       \rho_{\Sigma,\nu},\rho_{\Pi,\nu},c_{\Sigma,\nu},c_{\Pi,\nu}).
       \]
 \item \(\operatorname{or}_{{\rm coh},\lambda}\) : les signes
       \(\epsilon_\nu\), les ordres des ports et les applications
       \(\Theta_{\nu,N}\).
 \item \(\operatorname{Carr}_{{\rm coh},\lambda}\) : les matrices
       \(\beta_N\), les cocycles \(c_N(v,u)\), les résidus et les quotients,
       conservés séparément.
 \item \(\operatorname{Mem}_{{\rm coh},\lambda}\) : le ledger hérité et le
       mot ordonné des raffinements, retenues, corrections et avancées
       \(\Gamma_9\), non seulement leur valeur agrégée.
 \item \(\partial_{{\rm coh},\lambda}\) :
       \(\kappa_{\nu,N}\), \(D_{\rm tot}\),
       \(\partial_{\rm Cone}\) et les frontières des cellules.
 \item \(\gamma_{{\rm coh},\lambda}\) : toutes les chaînes composables
       \(\nu_0\to\cdots\to\nu_q\) et leurs continuations terminales.
 \item \(\operatorname{Msect}_{{\rm coh},\lambda}\) : toutes les cellules,
       les chemins terminaux, les extensions \(\Sigma\) et les corrections
       \(\operatorname{Corr}\), sans sélecteur.
 \item \(\operatorname{Surv}_{{\rm coh},\lambda}\) : le système des fibres
       \(\mathscr F_\lambda(a)\) et des restrictions surjectives.
 \item \(\operatorname{Term}_{{\rm coh},\lambda}\) :
       \(\operatorname{Cone}(I-U_{\Gamma_9,N})\), avec \(z\), \(Uz\)
       et leur différence.
 \item \(\mathcal L_{{\rm coh},\lambda}\) :
       \((z,[z],Uz,[(0,z)])\) et toutes ses compatibilités de réduction.
\end{enumerate}

Pour \(\lambda'\geq\lambda\), le transport coordonné satisfait
\begin{equation}\label{eq:pdfv2-coh-d-naturality}
 u_{\lambda',\lambda}^{\rm coh}
 D_{{\rm coh},\lambda'}(x_{\lambda'})
 =D_{{\rm coh},\lambda}
  (\rho_{\lambda',\lambda}x_{\lambda'}).
\end{equation}

\section{Graphe de restriction et objet dépendant}
\label{sec:pdfv2-coh-graphs}

La restriction totale est le graphe
\begin{equation}\label{eq:pdfv2-coh-g-graph}
 \mathcal G_{\rm coh}:=\operatorname{Graph}(D_{\rm coh}),
 \qquad
 \operatorname{Res}_{\rm coh}(x):=(x,D_{\rm coh}x).
\end{equation}
Ainsi,
\begin{equation}\label{eq:pdfv2-coh-res-inverse}
 \pi_1\operatorname{Res}_{\rm coh}=I,
 \qquad
 \operatorname{Res}_{\rm coh}\pi_1=I_{\mathcal G_{\rm coh}}.
\end{equation}

L'extracteur complet est
\begin{equation}\label{eq:pdfv2-coh-chi}
 \chi_{\rm coh}(x):=
 \bigl(
 \{C_{\nu,N}\},
 \{z_\lambda,[z_\lambda]\},
 \{\mathscr F_\lambda(e_\infty x)\},
 \{\Sigma_{\lambda',\lambda}\},
 \{\operatorname{Corr}_{\lambda',\lambda}\}
 \bigr).
\end{equation}
On définit
\begin{equation}\label{eq:pdfv2-coh-xchi}
 X_{\chi,{\rm coh}}
 :=\{(\chi_{\rm coh}(x),x,D_{\rm coh}x):
 x\in\mathcal C_{\rm cont}^{\rm coh}\}
\end{equation}
et
\begin{equation}\label{eq:pdfv2-coh-prx}
 \operatorname{pr}_{X,{\rm coh}}(x,D_{\rm coh}x)
 :=(\chi_{\rm coh}(x),x,D_{\rm coh}x).
\end{equation}
La projection sur \((x,D_{\rm coh}x)\) est inverse sur cet objet
dépendant.

\section{Assembleur et publicateur internes}
\label{sec:pdfv2-coh-assembler-publisher}

L'espace ambiant \(\mathscr M_{\rm coh}^{\rm amb}\) se compose de systèmes
compatibles de suites exactes de cellules, de catégories, de foncteurs, de modules de
chemins, de totalisateurs, de terminaux, de lecteurs, de fibres et de multisections.
L'assembleur brut est
\begin{equation}\label{eq:pdfv2-coh-assembler}
 \begin{aligned}
 a_{\rm coh}:X_{\chi,{\rm coh}}&\longrightarrow
 \mathscr M_{\rm coh}^{\rm amb},\\
 a_{\rm coh}(\chi(x),x,Dx)
 &:={}
 \bigl(
 \{C_{\nu,N},\kappa_{\nu,N},\delta_{\nu,N}\},
 \{J_d,\Gamma,W,K\},\\
 &\hspace{31mm}
 \{U_{\Gamma_9,N},\mathcal T_N\},
 \{e_\lambda,\mathscr F_\lambda,\Sigma,\operatorname{Corr}\}
 \bigr)_x.
 \end{aligned}
\end{equation}
La macrostructure est un autre graphe :
\begin{equation}\label{eq:pdfv2-coh-macro-graph}
 \mathfrak M_{\rm coh}:=\operatorname{Graph}(a_{\rm coh}),
 \qquad
 \mathcal A_{\rm coh}(z):=(z,a_{\rm coh}z).
\end{equation}

Soit \(\mathscr Y_{\rm coh}^{\rm int}\) le type de certificats qui contient
\begin{equation}\label{eq:pdfv2-coh-certificate-list}
 \begin{gathered}
 \text{la sucesión exacta \eqref{eq:pdfv2-coh-exact-sequence}},\\
 \kappa_{b,a}\tau=-P\kappa_{a,b},
 \qquad D_{\rm tot}^2=0,
 \qquad\partial_{\rm Cone}^2=0,\\
 rU_{\Gamma_9}=U_{\Gamma_9}r,
 \qquad
 e_\lambda\rho_{\lambda',\lambda}
 =r_{\lambda',\lambda}e_{\lambda'},\\
 \rho_{\lambda',\lambda}:\mathscr F_{\lambda'}(a)
 \twoheadrightarrow\mathscr F_\lambda(a),
 \qquad
 \Sigma_{\lambda',\lambda}(s_\lambda;a)\neq\varnothing,
 \end{gathered}
\end{equation}
avec les cycles bruts, leurs classes et leurs terminaux. Le publicateur
\begin{equation}\label{eq:pdfv2-coh-publisher}
 p_{\rm coh}^{\rm raw}:\mathfrak M_{\rm coh}
 \longrightarrow\mathscr Y_{\rm coh}^{\rm int}
\end{equation}
produit ce certificat sans supprimer la macrostructure. Enfin,
\begin{equation}\label{eq:pdfv2-coh-output-graph}
 \mathcal C_{\rm coh}^{\rm HMT}
 :=\operatorname{Graph}(p_{\rm coh}^{\rm raw}),
 \qquad
 \mathcal P_{\rm coh}(m):=(m,p_{\rm coh}^{\rm raw}m).
\end{equation}

\section{Théorème interne de génération cohérente}
\label{sec:pdfv2-coh-theorem}

\begin{theorem}[Génération interne de la macrostructure de cohérence]
\label{thm:pdfv2-coh-internal-generation}
Sur le domaine explicite \(\mathcal C_{\rm cont}^{\rm coh}\), la chaîne
\[
 \mathcal C_{\rm cont}^{\rm coh}
 \xrightarrow{\operatorname{Res}_{\rm coh}}
 \mathcal G_{\rm coh}
 \xrightarrow{\operatorname{pr}_{X,{\rm coh}}}
 X_{\chi,{\rm coh}}
 \xrightarrow{\mathcal A_{\rm coh}}
 \mathfrak M_{\rm coh}
 \xrightarrow{\mathcal P_{\rm coh}}
 \mathcal C_{\rm coh}^{\rm HMT}
\]
construit naturellement les cellules, leurs cônes, le diagramme d'histoires, le
bar bilatéral, la totalisation, le terminal nonadique et les fibres
survivantes. Elle conserve les treize coordonnées et toutes les corrections comme
multisections. Chaque étape admet, sur son image, une projection qui récupère
l'étape précédente.
\end{theorem}

\begin{proof}
La \cref{prop:pdfv2-coh-cell-exact} construit chaque cellule sur \(A_N\), avec
noyau diagonal et défaut \((a_\nu-1)(b_\nu-1)\). L'identité
\eqref{eq:pdfv2-coh-orientation} fait de l'échange orienté une application de
complexes involutive. Les formules
\eqref{eq:pdfv2-coh-gamma-one}--\eqref{eq:pdfv2-coh-gamma-zero} prouvent que
chaque raffinement est une application de complexes, tandis que
\eqref{eq:pdfv2-coh-kappa-reduction} prouve la compatibilité en précision.

La précomposition des chemins rend \(W\) contravariant. Les identités
fonctorielles de \(W\) et \(\Gamma\) impliquent
\(d_{\rm bar}^2=0\) ; leur commutation avec la différentielle cellulaire donne
\(D_{\rm tot}^2=0\). Ainsi, \(K_{d,N}\) est construit par les formules bar,
non par une invocation nominale du cofin.

L'extension \(\Gamma_9\) augmente la profondeur. Elle n'induit donc
l'endomorphisme \(U_{\Gamma_9,N}\) qu'après le passage à la colimite
\(K_{\infty,N}\). Le cône
\(\operatorname{Cone}(I-U_{\Gamma_9,N})\) conserve alors la différence
\([h]-[\Gamma_9h]\) et n'identifie pas la répétition de phase à la répétition
d'état.

Le cycle \(z_{d,N}(x)\) se forme avec toutes les cellules, tous les signes, chemins et
incidences de l'histoire. Troncature et réduction commutent avec le lecteur. Si
\(s_\lambda\in\mathscr F_\lambda(a)\), l'histoire complète qui apparaît dans
la définition même de la fibre fournit un antécédent à tout niveau
\(\lambda'\geq\lambda\) ; cela prouve la surjectivité
\eqref{eq:pdfv2-coh-fiber-surjection}. L'objet conserve l'ensemble entier
\(\Sigma_{\lambda',\lambda}\), non un choix. Les différences entre
relèvements sont enregistrées par
\eqref{eq:pdfv2-coh-correction-relation}, avec support, orientation, retenue,
mémoire et terminal explicites.

Enfin, \(\mathcal G_{\rm coh}\), \(\mathfrak M_{\rm coh}\) et
\(\mathcal C_{\rm coh}^{\rm HMT}\) sont des graphes, tandis que
\(X_{\chi,{\rm coh}}\) conserve comme coordonnées \(x\) et \(D_{\rm coh}x\).
La première projection de chaque objet récupère le précédent. Par conséquent, la
composition complète ne remplace pas l'histoire par une classe, la multisection par un
sélecteur ni le terminal par la phase visible.
\end{proof}

\section{Provenance probatoire et falsificateurs}
\label{sec:pdfv2-coh-provenance-falsifiers}

\paragraph{Résultat récupéré.}
Appartiennent à l'architecture déjà construite la cellule \(\kappa_{a,b}\), son
noyau diagonal et son défaut, l'orientation
\(\kappa_{b,a}\tau=-P\kappa_{a,b}\), la retenue APP, le cône de la cellule, la
catégorie d'histoires, la totalisation par cofin, l'action \(\Gamma_9\), le
terminal \(\operatorname{Cone}(I-U_{\Gamma_9})\) et l'architecture
d'extension survivante.

\paragraph{Formalisation nouvelle.}
Sont des formalisations explicites la suite exacte
\eqref{eq:pdfv2-coh-exact-sequence} sur \(A_N\), la séparation entre
profondeur et précision, les faces complètes du bar, la formation du
terminal seulement après la colimite, la mise en paquet des treize coordonnées,
les corrections comme relation sans sélecteur et les quatre objets dépendants
ou graphes qui empêchent la perte de généalogie.

\begin{criterion}[Falsificateurs internes de la branche \(\mathrm{coh}\)]
\label{crit:pdfv2-coh-falsifiers}
La construction échoue si l'un quelconque des faits suivants se produit :
\begin{enumerate}
 \item la suite
       \eqref{eq:pdfv2-coh-exact-sequence} n'est pas exacte ;
 \item \(\kappa_{b,a}\tau=-P\kappa_{a,b}\) échoue ;
 \item un raffinement ne satisfait pas
       \(\kappa_{\nu'}\Gamma(u)_1=\Gamma(u)_0\kappa_\nu\) ;
 \item la retenue de composition ne satisfait pas
       \eqref{eq:pdfv2-coh-carry-cocycle} ;
 \item \(D_{\rm tot}^2\neq0\);
 \item la réduction, le lecteur ou le terminal ne commutent pas avec \(U_{\Gamma_9}\) ;
 \item on forme \(\operatorname{Cone}(I-U_{\Gamma_9})\) en identifiant d'abord
       les profondeurs \(d\) et \(d+9\) ;
 \item une fibre déclarée survivante est vide ou une restriction n'est pas
       surjective ;
 \item on remplace \(\Sigma_{\lambda',\lambda}\) par un unique
       relèvement ;
 \item on publie seulement \([z]\) et l'on supprime \(z\), le chemin, la retenue, la mémoire ou
       le terminal ;
 \item l'une quelconque des treize coordonnées est omise.
\end{enumerate}
Dans le dernier cas, le diagnostic contraignant est : \emph{objet substitué ;
conclusion non transférable}.
\end{criterion}

\begin{warningbox}
Ce chapitre conclut uniquement la construction interne de
\(\mathcal C_{\rm coh}^{\rm HMT}\). Il n'introduit pas d'espace classique, de
classe classique ni de publication classique. Tout entrelaceur ultérieur
appartient à un autre document et ne gouverne pas l'existence de cette
macrostructure.
\end{warningbox}
\fi
